\par\Needspace{6\baselineskip}\medskip
\begin{center}\hypertarget{native-246-article175}{}\pdfbookmark[2]{175.}{article-175}175.\end{center}
Zu diesem Zweck gehen wir von unserer besonderen Aufgabe zu einer sehr allgemeinen Untersuchung über, die sich bei jeder. Anwendung der Rechnung auf die Naturphilosophie als äußerst fruchtbar erweist. Es seien $V$, $V'$, $V''$ usw. Funktionen der Unbekannten $p$, $q$, $r$, $s$ usw., $\mu$ sei die Anzahl dieser Funktionen und $\nu$ die Anzahl der Unbekannten. Wir nehmen an, dass sich durch unmittelbare Beobachtungen die folgenden Funktionswerte ergeben haben: $V=M$, $V'=M'$, $V''=M''$ usw. Im Allgemeinen ist also die Ermittlung der Werte der Unbekannten eine unbestimmte, bestimmte oder überbestimmte Aufgabe, je nach der geltenden Beziehung $\mu<\nu$, $\mu=\nu$ oder $\mu>\nu$\footnote{Wären im dritten Fall die Funktionen $V$, $V'$, $V''$ usw. so beschaffen, dass sich $\mu+1-\nu$ oder mehr von ihnen als Funktionen der übrigen ansehen ließen, so wäre die Aufgabe hinsichtlich dieser Funktionen immer noch überbestimmt, hinsichtlich der Größen $p$, $q$, $r$, $s$ usw. aber unbestimmt: Ihre Werte ließen sich nämlich selbst dann nicht bestimmen, wenn die Werte der Funktionen $V$, $V'$, $V''$ usw. absolut genau gegeben wären: Diesen Fall schließen wir jedoch aus unserer Untersuchung aus.}.
Hier ist nur vom letzten Fall die Rede, in dem sich offenbar alle Beobachtungen nur dann exakt darstellen ließen, wenn sie sämtlich völlig fehlerfrei wären. Da dies in der Wirklichkeit nicht vorkommt, muss jedes System von Werten der Unbekannten $p$, $q$, $r$, $s$ usw. als möglich gelten, aus dem sich Werte der \SourceLeaf{247}Funktionen $M-V$, $M'-V'$, $M''-V''$ usw. ergeben, die nicht größer sind als die Grenzen der Fehler, die bei diesen Beobachtungen auftreten konnten: Das ist jedoch keineswegs so zu verstehen, als ob alle diese möglichen Systeme den gleichen Grad an Wahrscheinlichkeit hätten.

Nehmen wir zunächst an, die Umstände seien bei allen Beobachtungen so gewesen, dass kein Grund zu der Vermutung besteht, eine sei weniger genau als eine andere, oder, anders ausgedrückt, dass gleich große Fehler bei jeder Beobachtung als gleich wahrscheinlich anzusehen sind. Die Wahrscheinlichkeit, die einem beliebigen Fehler $\Delta$ zukommt, wird folglich durch eine Funktion von $\Delta$ selbst ausgedrückt, die wir mit $\varphi\Delta$ bezeichnen. Auch wenn sich diese Funktion nicht genau angeben lässt, können wir doch sagen, dass ihr Wert für $\Delta=0$ am größten sein muss, dass er für betragsgleiche Werte von $\Delta$ mit entgegengesetztem Vorzeichen in der Regel gleich ist und schließlich verschwindet, wenn man für $\Delta$ den größten Fehler oder einen größeren Wert einsetzt. Streng genommen ist $\varphi\Delta$ somit den unstetigen Funktionen zuzurechnen, und wenn wir uns für praktische Zwecke erlauben, sie durch eine analytische Funktion zu ersetzen, so muss diese so beschaffen sein, dass sie auf beiden Seiten von $\Delta=0$ gewissermaßen asymptotisch gegen $0$ konvergiert und daher jenseits dieser Grenze als tatsächlich verschwindend betrachtet werden kann. Ferner ist die Wahrscheinlichkeit, dass der Fehler zwischen den Grenzen $\Delta$ und $\Delta+\mathrm{d}\Delta$ liegt, deren Abstand die unendlich kleine Differenz $\mathrm{d}\Delta$ beträgt, durch $\varphi\Delta\, .\,\mathrm{d}\Delta$ auszudrücken; im Allgemeinen ergibt sich daher die Wahrscheinlichkeit, dass der Fehler zwischen $D$ und $D'$ liegt, aus dem Integral $\int\varphi\Delta\, .\,\mathrm{d}\Delta$, genommen von $\Delta=D$ bis $\Delta=D'$. Dieses Integral, genommen vom äußersten negativen Wert von $\Delta$ bis zum größten positiven Wert, oder allgemeiner von $\Delta=-\infty$ bis $\Delta=+\infty$, muss notwendigerweise $=1$ sein.

Nehmen wir also an, ein bestimmtes System von Werten der Größen $p$, $q$, $r$, $s$ usw. liege vor. Die Wahrscheinlichkeit, dass sich für $V$ aus der Beobachtung der Wert $M$ ergibt, wird durch $\varphi(M-V)$ ausgedrückt, nachdem man in $V$ für $p$, $q$, $r$, $s$ usw. die jeweiligen Werte eingesetzt hat; ebenso drücken $\varphi(M'-V')$, $\varphi(M''-V'')$ usw. die Wahrscheinlichkeiten aus, dass die Beobachtungen für die Funktionen $V'$, $V''$ usw. die Werte $M'$, $M''$ usw. ergeben. Da sich somit alle Beobachtungen als voneinander unabhängige Ereignisse ansehen lassen, drückt das Produkt
\[
\varphi(M-V)\, .\,\varphi(M'-V')\, .\,\varphi(M''-V'')\, .\,\text{usw.}=\Omega
\]
die Erwartung oder Wahrscheinlichkeit aus, dass sich alle diese Werte gemeinsam aus den Beobachtungen ergeben.
